In this section we prove our main result.  In the previous section we showed how to compute the girth of an element in a free group.  The Main Theorem applies to curves in a subsurface of the boundary of a handlebody $H$.

\begin{theorem}[\textbf{Main Theorem}]
\label{subsurface theorem}
Let $H$ be a handlebody.  Let $S$ be a compact surface in $\partial H$.  The minimum girth in $H$ over all essential closed curves in $S$ can be achieved by a simple closed curve in $S$.
\end{theorem} 

We begin with some definitions and lemmas.  Let $H$ be a handlebody, and let $\gamma$ be a disjoint collection of simple closed curves in $\partial H$.  Let $\dd$ be a disc system for $H$.  Define the \emph{geometric complexity} of $\dd$ with respect to $\gamma$ to be the intersection number $|\dd\cap\gamma|$.  This parallels the notion of homotopic complexity for a word $w$ in a free group that we presented in Section~\ref{computecurve}.  We will sometimes say just \emph{complexity}.  Define a disc system to be \emph{taut} if it minimizes the complexity in its isotopy class.

Let $\dd$ be a taut disc system with respect to $\gamma$.  Recall that $\dpm$ denotes the $2g$ oriented discs coming from $\dd$, where $g$ is the genus of $H$.  Define a \emph{half-disc} based at an oriented disc $D$ in $\dpm$ to be an arc $K$ in $\bd H$ with endpoints in $\bd D-\gamma$ and interior embedded essentially in $\partial H-\bd \dd$.  By essentially we mean that, inside $\bd H$, $K$ cannot be isotoped $\rel \bd K$ into $\bd \dd$.  We also require that $K$ enter and leave $D$ from the side corresponding to its orientation.  See Figure~\ref{halfdisc_shortcut} for two examples of half-discs.  Observe that, inside $H$, a half-disc $K$ can be isotoped $\rel \bd K$ into $D$.

\begin{figure}
% a shortcut half disc and a non-shortcut half-disc
\epsfig{file=figures/halfdisc_shortcut.eps, height=2in}
\caption{Both pictures illustrate a half-disc $K$ based at a shaded disc $D$.  The right one is a short-cut half-disc, and the left one is not.}
\label{halfdisc_shortcut}
\end{figure}
 
Let $K$ be a half-disc based at a disc $D$ in $\dpm$.  The endpoints of $K$ divide $\bd D$ into two components $\bd D_B$ and $\bd D_W$.  We say that $K$ is a \emph{shortcut} half-disc if $|K\cap\gamma| < \min(|D_B\cap\gamma|,|D_W\cap\gamma|)$.  See Figure~\ref{halfdisc_shortcut} for an example.  By the following lemma, if $\dd$ has a shortcut half-disc, we can construct a new disc system with smaller complexity.  Lemma~\ref{shortcut} is the geometric analogue of Lemma~\ref{curve shortcut}.

\begin{lemma}[cf. Lemma~\ref{curve shortcut}]
\label{shortcut}
Let $H$ be a handlebody, and let $\gamma$ be a disjoint collection of simple closed curves in $\partial H$.  Let $\dd$ be a taut disc system with respect to $\gamma$.  If there is a shortcut half-disc with respect to $\gamma$, then there is a disc system with smaller complexity than that of $\dd$.
\end{lemma}

\begin{proof}
Let $K$ be a shortcut half-disc based at an oriented disc 
$D$ in $\dpm$.  Let $\bd D_B$ and $\bd D_W$ be the two arc components of $\partial D-\partial K$.
Since $K$ is a shortcut half-disc with respect to $\gamma$, 
we know 
$|K\cap\gamma|<\min(|\bd D_B\cap\gamma|,|\bd D_W\cap\gamma|)$.
Boundary-compress $D$ along $K$ to create two essential 
discs in $H$ that are disjoint from $\dd$.  See Figure~\ref{halfdisc_compress}.
By the same argument used in the proof of Lemma~\ref{curve shortcut}, replacing $D$ in $\dd$ with one of these two discs, say $D^\prime$, yields a disc system.  Replace $D$ with $D^\prime$ to form a new disc system $\dd^\prime$.  Say $D^\prime$ was the disc corresponding to the arc $\bd D_B$.  Since
$$|D^\prime\cap\gamma|=|D_B\cap\gamma|+|K\cap\gamma|<
|D_B\cap\gamma|+|D_W\cap\gamma|=|D\cap\gamma|,$$ 
the disc system $\dd^\prime$ has a complexity smaller than the complexity of $\dd$.
\end{proof}

\begin{figure}
% the two discs from a half-disc
\epsfig{file=figures/halfdisc_compress.eps, height=2.25in}
\caption{The left picture shows a compressing half-disc $K$ based at a shaded disc $D$.  The right picture shows the disc $D$ after boundary-compressing along $K$.}
\label{halfdisc_compress}
\end{figure}

We say a disc system $\dd$ is \emph{locally minimal} with respect to $\gamma$ if there are no shortcut half-discs.  Observe that the proof of Lemma~\ref{shortcut} is constructive, so it gives an algorithm to construct a locally minimal disc system, given any collection of curves $\gamma$ in $\bd H$.

Define a \emph{compressing half-disc} to be a half-disc $K$ with $|K\cap\gamma|=0$.  See Figure~\ref{halfdisc_compress}.  We say a disc system $\dd$ is \emph{compressible} if $\dd$ has a compressing half-disc.  By the following lemma, if a curve collection $\gamma$ has a compressible locally minimal disc system, then the surface $\bd H-\gamma$ is compressible in $H$.

\begin{lemma}
\label{no waves}
Let $H$ be a handlebody, and let $\gamma$ be a disjoint collection
of simple closed curves in $\partial H$.
Let $\dd$ be a locally minimal disc system with respect to $\gamma$.  
If $\dd$ is compressible, then $H$ contains an essential disc disjoint
from $\dd$ and $\gamma$.
\end{lemma}

\begin{proof}
Let $K$ be a compressing half-disc for $\dd$ based at a disc $D$ in $\dpm$, so that $|K\cap\gamma|=0$.  Let $\bd D_B$ and $\bd D_W$ denote the two components of $\bd D-\bd K$.  Since $\dd$ is locally minimal, the half-disc $K$ is not a shortcut half-disc.  Thus $|K\cap\gamma|\geq\min(|\bd D_B\cap\gamma|,|\bd D_W\cap\gamma|)$.  In other words, at least one of the two arcs, say $\bd D_B$, satisfies $|\bd D_B\cap\gamma|=0$.  Boundary-compress $D$ along $K$ to create two essential discs in $H$ that are disjoint from $\dd$.  The disc corresponding to $\bd D_B$ can be made disjoint from $\gamma$ since neither $K$ nor $\bd D_B$ intersected $\gamma$.
\end{proof}

Let $S$ be a compact surface.  Let $\gamma$ be an essential closed curve in $S$.  We say that $\gamma$ \emph{fills $S$} if $\gamma$ cannot be homotoped into a subsurface $S^\prime$ of $S$ with $\ec{S^\prime}<\ec{S}$.  See Figure~\ref{fillingcurve}.  Let $S^\prime$ be a subsurface of $S$.  We say that $S^\prime$ is \emph{incompressible} in $S$ if $\pi_1(S^\prime)$ injects into $\pi_1(S)$.  By the following lemma, any curve in a surface $S$ can be homotoped to fill some incompressible subsurface in $S$.

\begin{figure}
% a curve that's not full
\epsfig{file=figures/fillingcurve.eps, height=1.5in}
\caption{The curve in the left picture fills the surface.  The curve on the right does not fill.  It can be homotoped into an annulus.}
\label{fillingcurve}
\end{figure}

% \begin{figure}
% a surface compressible in another surface
% \epsfig{file=images/09.eps, height=2.25in}
% \caption{The surface $S^\prime$ shown here is compressible in $S$ 
% because it has a hole that can be compressed, i.e. filled in.}
% \label{09}
% \end{figure}

\begin{lemma}
\label{full}
Let $S$ be a compact surface with boundary.  Let $\gamma$ be an 
essential closed curve in $S$.  Then there is an incompressible subsurface $S^\prime$ of $S$ and a curve $\gamma^\prime$ such that $\gamma^{\prime}$ both fills $S^{\prime}$ and is homotopic to $\gamma$ in $S$.
\end{lemma}

\begin{proof}
The proof is by downward induction on $\ec{S}$.  If $\ec{S}=0$, then $\gamma$ fills $S$.  Assume that $\ec{S}>0$.  If $\gamma$ does not fill $S$,
then homotope $\gamma$ into a compact subsurface $S^\prime$ of $S$ satisfying $\ec{S^\prime}<\ec{S}$.  If $S^\prime$ is compressible, then compress $S^\prime$.  The inequality $\ec{S^\prime}<\ec{S}$ still holds.  Continue compressing $S^\prime$ until it  is incompressible in $S$.  Rename $S^\prime$ as $S$, and repeat the induction.
\end{proof}

Let $H$ be a handlebody of genus $g$, and let $S$ be a compact subsurface of  $\partial H$.  Let $\dd$ be a taut disc system with respect to $\partial S$.  We associate the following diagram to the triple $(H,S,\dd)$.  Take the complement in $H$ of an open regular neighborhood of $\dd$.  This defines a closed ball in $H$.  The boundary of this ball is a sphere $\m$.  The sphere $\m$ contains $2g$ distinguished discs that correspond to the discs of $\dpm$.  We call these discs the \emph{countries} of $\m$.  The surface $S$ intersected with $\m$ comprises several polygonal regions.  We call these regions the \emph{roads}.  The diagram of $\m$, together with its countries and roads, is the \emph{map} associated to $(H,S,\dd)$.  See Figure~\ref{map} for a picture.

Maps can be illustrated in the plane as follows.  Let $\m$ be a map.  View $\m$ as $\R^2\cup\infty$, the plane with a point at infinity.  Choose a country of $\m$.  Put that country around the point at infinity.  Now draw the diagram $\m$ as a disc with the other $2g-1$ countries inside.  
\begin{figure}
% the associated diagram
\epsfig{file=figures/map.eps, height=1.75in}
\caption{Here is a map where the handlebody has genus two and the $\dd$-interface has three white countries and one black country.
}
\label{map}
\end{figure}

A $\dd$-interface adds color to a map.  Let $\m$ be a map, and let $\dint{a}$ be a $\dd$-interface.  The $\dd$-interface  $\dint{a}$ partitions $\dpm$ into black, white, and gray discs.  Color the corresponding countries of $\m$ accordingly.  Again, see Figure~\ref{map}.

Let $R$ be a road of the map $\m$.  We say the road $R$ \emph{segregates} 
if (1) it borders a black country and a white country but no gray country, and (2) a single embedded arc can separate its black borders from its white borders.  We call such an arc a \emph{line of segregation}.  See Figure~\ref{segregation}.

\begin{figure}
% segregated and non-segregated roads
\epsfig{file=figures/segregated.eps, height=1.5in}
\caption{Here are three illustrations of roads and the countries they border.  The road on the left segregates.  The dotted line shown is a line of segregation.  The roads on the right do not segregate.
}
\label{segregation}
\end{figure}

The following lemma relates nearly all of the concepts just introduced.

\begin{lemma}
\label{segregated regions}
Let $H$ be a handlebody, and let $S$ be a compact subsurface of $\partial H$.
Let $\dd$ be a taut, incompressible disc system with respect to $\bd S$.  Let $\dint{a}$ be a $\dd$-interface.  Let $\m$ be the associated map.  Let $\gamma$ be a closed curve in $S$, and suppose that $\gamma\wedge\dint{a}=n$.  If $\gamma$ fills $S$, then $\m$ has at most $n$ segregating roads.
\end{lemma}

\begin{proof}
Lift $\dint{a}$ to the universal cover $\hh$, and call the lift $\alpha$.  Exactly $n$ lifts of $\gamma$ cross $\alpha$.  Let $\db$ and $\dw$ denote the black and white discs of $\dpm$ associated with $\alpha$.

Since $\dd$ is incompressible with respect to $\bd S$, there are no compressing half-discs in $\bd H$.  Thus every curve in $S$ is isotopic to a geodesic with respect to $\dd$.  Isotope $\gamma$ to be a geodesic with respect to $\dd$.

View $\gamma$ in $\m$, so that $\gamma$ comprises a collection of arcs inside $\m$'s roads.  By Lemma~\ref{intersection handlebody}, every arc of $\gamma$ that connects a white country to a black country corresponds to a lift of $\gamma$ that crosses $\alpha$.

We proceed by contradiction.  Assume that $\m$ has more than $n$ segregated roads.  Then, in at least one of these roads, no arc of $\gamma$ passes from a white country to a black country.  Let $R$ be such a road.  Choose a line $\ell$ of segregation for $R$.  Isotope $\gamma$ so that it does not intersect $\ell$.  Then $S-\ell$ is a surface containing $\gamma$ and satisfying $\ec{S-\ell}<\ec{S}$.  This contradicts that $\gamma$ is full.
\end{proof}

Our final lemma is a topological result about arcs in surfaces.

\begin{lemma}
\label{arcs in S}
Let $S$ be a compact surface with boundary.  Let $n$ be a natural number such that $n\geq0$.  Let $K$ be a collection of $\ec{S}+n+1$ arcs in $S$ that are disjoint, essential, and embedded $\rel \bd S$.  If $\chi(S)\leq -1$ and $n\geq 2$, then $S$ contains an essential simple closed curve that intersects $K$ no more than $n$ times.  Otherwise, $S$ contains an essential simple closed curve that intersects $K$ no more than $n+1$ times.
\end{lemma}

\begin{proof}
It is obvious when $\chi(S)=0$ that $S$ contains an essential simple closed curve that intersects $K$ exactly $n+1$ times.  So assume $\chi(S)\leq -1$.

Assume the complement of $K$ has a component that is not simply-connected.  Then there is an essential simple closed curve in $S$ that intersects $K$ zero times.  We are done in this case, so assume otherwise.

Choose a minimal collection $K_{0}$ of arcs that cut $S$ into a disc.  Any such collection contains $\ec{S}+1$ arcs.  Cut $S$ along $K_{0}$ to form a polygonal region $P$.  The polygon $P$ contains $n$ arcs leftover in its interior.  If $n=0$, there are essential curves in $S$ that intersect $K$ one time, and we would be done.  So assume that $n\geq 1$.

Let $K_{1}$ denote the $n$ interior arcs in $P$.  Equivalently, $K_{1}$ equals $K-K_{0}$.  Each arc in $K_{0}$ corresponds to a pair of sides in $\partial P$.  Consider such a pair of sides $\beta$ in $P$.  Each arc $\kappa$ in $K_{1}$ either separates or does not separate $\beta$.  An arc $\kappa$ is said to \emph{separate} $\beta$ if and only if $\beta$ cannot be connected by a path in $P$ that is disjoint from $\kappa$.

Assume there is an arc $\kappa$ in $K_{1}$ that does not separate some pair $\beta$.  Then $\beta$ can be connected by a path in $P$ that intersects no more than $n-1$ arcs in $K_{1}$.  In other words, the surface $S$ contains an essential simple closed curve that intersects no more than $n$ arcs in $K$.  In this case, we are done.  So assume otherwise, namely that all arcs in $K_{1}$ separate all pairs $\beta$.

For every pair $\beta$ in $P$, all paths connecting $\beta$ in $P$ intersect all $n$ arcs in $K_{1}$.  This implies that the $n$ interior arcs $K_{1}$ are parallel, as shown in Figure~\ref{polygonarcs}.  Since $S$ is not annulus, we can find an essential simple closed curve in $S$ that intersects just two arcs in $K_{0}$ and avoids all $n$ arcs in $K_{1}$.  Again, see Figure~\ref{polygonarcs} for an illustration.  Since $S$ contains a curve intersecting $K$ exactly two times, the proof is complete for both $n=1$ and $n\geq 2$.
\end{proof}

\begin{figure}
% the polygon P
\epsfig{file=figures/surfacelemma.eps, height=2.25in}
\caption{Here is an example surface $S$ for the proof of Lemma~\ref{arcs in S}.  We have $\chi(S)=-2$ and $n=3$.  Cut $S$ along $\ec{S}+1$ arcs to form the polygonal disc shown above.  Arrows indicate the gluing instructions.  The dotted line represents a curve intersecting $K$ just twice.
}
\label{polygonarcs}
\end{figure}

We now prove the Main Theorem.  We state and prove 
Lemmas~\ref{main lemma easy} and~\ref{main lemma corollary} 
in the next section, Section~\ref{mainlemma}.

\begin{proof}[Proof of Theorem~\ref{subsurface theorem}]
Assume without loss of generality that $S$ is incompressible in $\bd H$.  Then no component of $\bd S$ bounds a disc in $\bd H$.

Let $\gamma$ be a curve in $S$ that minimizes the girth in $H$ over all essential closed curves in $S$.  Let $n$ be the girth of $\gamma$.  Apply Lemma~\ref{full}, and homotope $\gamma$ into an incompressible subsurface $S^\prime$ of $S$ so that $\gamma$ fills $S^\prime$.

It suffices to prove the theorem for $S^\prime$.  Give $S^\prime$ the name $S$ for convenience.  If $S$ is an annulus, then we are done by Proposition~\ref{power}: a core curve for the annulus minimizes the girth.  So assume $S$ is not an annulus.  In particular, we can also assume $\genus(H)\geq2$.

The curve $\gamma$ fills the surface $S$.  We will construct an essential disc $D$ in $H$ and an essential simple closed curve $\gamma^\prime$ in $S$ so that $\gamma^\prime$ intersects $D$ in $n$ points.

Using Lemma~\ref{shortcut}, find a locally minimal disc system $\dd$ with respect to $\bd S$.  Then by Lemma~\ref{no waves}, either $\dd$ is incompressible with respect to $\bd S$, or there is a disc $D$ in $H-\dd$ disjoint from $\partial S$.  In the latter case, $\bd S$ contains simple closed curves of girth zero, and we are done.  So assume that $\dd$ is incompressible.

Let $B$ be a girth-realizing disc for $\gamma$, so that $\gamma\wedge B=n$.  Apply Lemma~\ref{subordination-lemma} to $B$.  There are two cases.

\textbf{Case 1:  There is a $\dd$-disc $\dint{b}$ subordinate to $B$.}
Then $\gamma\wedge\dint{b}\leq n$.  By Lemma~\ref{segregated regions}, the $\dd$-disc $\dint{b}$ has no more than $n$ segregated roads.  Apply Lemma~\ref{main lemma easy} to obtain a disc $D$ in $H-\dd$ that intersects $S$ in no more than $\ec{S}+n+1$ arcs.  Now proceed to the conclusion following Case 2.

\textbf{Case 2:  There are two $\dd$-half-discs independently subordinate
to $A$.}  Let $\dint{c}$ be the $\dd$-half-disc that $\gamma$ homotopically 
intersects less.  Let $m$ equal $\gamma\wedge\dint{c}$.  Then $m\leq\lfloor\frac{n}{2}\rfloor$.  By Lemma~\ref{segregated regions}, the $\dd$-half-disc $\dint{c}$ has no more than $m$ segregated roads.
Apply Lemma~\ref{main lemma corollary} to obtain a disc $D$ in $H-\dd$
that intersects $S$ in no more than $\ec{S}+2m+1\leq\ec{S}+n+1$ arcs.
Proceed now to the conclusion below.

\textbf{Conclusion of Cases.}
So far we have constructed a disc $D$ in $H$ that intersects $S$ in no more than $\ec{S}+n+1$ arcs.  Since $S$ is not an annulus, we have $\ec{S}\leq-1$.  Apply Lemma~\ref{arcs in S} to $S$.  There are two possible conclusions.

If $n=0$, Lemma~\ref{arcs in S} gives us an essential simple closed curve $\gamma^{\prime}$ in $S$ that intersects $D$ no more than once.  If necessary, apply Lemma~\ref{geometric girth one} to construct a disc that intersects $\gamma^\prime$ zero times.  This completes the case of $n=0$.

Assume next that $n\geq1$.  Since $n$ is the girth of a curve, we have by Lemma~\ref{group girth one} that $n\neq1$.  Hence $n\geq2$.  In this case, Lemma~\ref{arcs in S} gives us an essential simple closed curve $\gamma^{\prime}$ in $S$ that intersects $D$ no more than $n$ times.  Thus we are done.
\end{proof}

